\section{Running of the Wilson coefficients}
\label{sec:rge}

The parameterization of Sec.~\ref{sec:parameterization} treats the Wilson coefficients as
constants. They are supplied to the generator at a fixed input scale, $\mu_0 = 173$~GeV,
and the matrix element uses that value at every phase-space point. In the SMEFT the
coefficients run, and a calculation that resums the associated logarithms evaluates them
at a scale characteristic of the process, which for \HH production grows with \mHH. Like
the validity statements of Sec.~\ref{sec:truncation}, this is a property of
Eq.~\eqref{eq:smeftpara_reorg} and not of any analysis selection.

Five of the nine coefficients evolve --- $c_{\PQt\PH}$, $c_{\PH G}$, $c_{\PQt G}$,
$c_{Q\PQt}^{(1)}$ and $c_{Q\PQt}^{(8)}$ --- while the purely bosonic $c_{\PH\Box}$,
$c_{\PH D}$ and $c_\PH$ do not. The subleading coefficients evolve multiplicatively from
their input values, so the operator mixings $c_{Q\PQt}^{(1,8)} \to c_{\PQt\PH}$ and
$c_{\PQt G}, c_{Q\PQt}^{(1,8)} \to c_{\PH G}$ vanish at a point where the subleading
coefficients do. The anomalous dimensions and their solution are those of
Ref.~\cite{Heinrich:2024rtg}.

The size of the omission is measured on samples generated with the running enabled at
$\sqrt{s}=13.6$~TeV, with $\mu_{\text{EFT}} = \mHH/2$. A benchmark with no running
coefficient is reproduced to $0.02\%$ in rate, which bounds any artefact of the
procedure. At $c_{\PQt\PH}=4$, the largest value of the benchmark set, the rate falls by
$2.9\%$ and $\langle\mHH\rangle$ by $1.5\%$, with a coherent tilt of the spectrum at
$5.8\sigma$ --- soft bins up and hard bins down, as expected when a coefficient runs down
with \mHH --- and the effect falls with $c_{\PQt\PH}$.

The coefficients delivered with this report are therefore defined at $\mu_0 = 173$~GeV
and the templates are built with them fixed there. Where the coefficients that run are
small the difference is below the statistics of the comparison; at the upper end of the
$c_{\PQt\PH}$ range it is a few per cent in rate and about one per cent in the mean of
the spectrum, comparable to the modelling uncertainty of Sec.~\ref{sec:uncertainties} and
with a definite shape rather than a normalization. An interpretation that quotes
coefficients at a different scale, or that combines these results with measurements which
do, should either evolve them explicitly or carry an uncertainty of that size. For the
four coefficients that do not run there is nothing to do.

Evolving them explicitly costs no new simulation. At leading logarithmic accuracy the
solution of the dimension-6 renormalization-group equations is linear in the Wilson
coefficients, $c_i(\mu) = U_{ij}(\mu,\mu_0)\,c_j(\mu_0)$, and a linear map composed with
the quadratic polynomial of Eq.~\eqref{eq:smeftpara_reorg} is again a quadratic
polynomial in the same monomials (Sec.~\ref{sec:rge_generation}); the morphing weights
are polynomials in those same monomials, so the basis delivers the running prediction at
the transformed point with the templates it already has.

The shape moves as well, and by more than the rate where several coefficients run
together: over the benchmark set of Sec.~\ref{sec:trunc_modelling}, extended with three
nine-dimensional points, the mean of the \mHH spectrum shifts by up to $9\%$, the
largest shift being at the point that isolates the gluonic and four-top channels. An
analysis sensitive to the spectrum therefore needs the kinematic description to follow
the running and not only the rate. The implementation of that map, its closure against
the running-enabled samples by both routes, and the benchmarks in which the four-fermion
mixings are active are going to be studied in a follow-up work.
